\chapter{Teleoperation and Data Recording}
\label{chap:touch}
\label{chap:collection}

Teleoperation is not only a convenient way to move the robot. It defines the
behavior that the policy will imitate. Controller lag, inconsistent gripper
commands, occlusion, and unnecessary corrective motions all become part of the
training target.

\section{A reusable teleoperation architecture}

A practical teleoperation system separates device input, target generation,
low-level control, and recording:

\begin{center}
\resizebox{\textwidth}{!}{%
\begin{tikzpicture}[
  node distance=7mm and 7mm,
  box/.style={draw=lirmmblue,rounded corners=1mm,thick,align=center,
              minimum height=10mm,minimum width=27mm,fill=softblue},
  arr/.style={-{Latex[length=2.2mm]},thick,lirmmblue}
]
  \node[box] (device) {Operator\\device};
  \node[box,right=of device] (ros) {ROS input\\messages};
  \node[box,right=of ros] (target) {Desired pose\\and gripper};
  \node[box,right=of target] (servo) {Robot-specific\\servo};
  \node[box,right=of servo] (robot) {Simulation or\\physical robot};
  \node[box,below=9mm of target] (record) {Synchronized\\recorder};
  \draw[arr] (device) -- (ros);
  \draw[arr] (ros) -- (target);
  \draw[arr] (target) -- (servo);
  \draw[arr] (servo) -- (robot);
  \draw[arr] (robot.south) |- (record.east);
  \draw[arr] (target.south) -- (record.north);
\end{tikzpicture}%
}
\end{center}

This separation makes the input device and learned-policy interface reusable.
Only the servo and state reader must change for another robot.

\begin{projectbox}[Implementation used during the internship]
The 3D Systems Touch publishes pose and buttons through a ROS~2 C++ node. A
Python listener converts incremental stylus motion into a Cartesian target.
Pinocchio computes forward kinematics and the Jacobian
\citep{carpentier2019pinocchio}. The same target logic feeds either MuJoCo or a
UR10 \filepath{speedj} backend.
\end{projectbox}

\section{Use relative motion, not device coordinates}

The physical coordinates of a desktop device generally have no useful absolute
relationship with the robot base. Use differential motion:

\begin{equation}
  \Delta p_{robot} = s_p M_p
  \left(p^{device}_{t}-p^{device}_{t-1}\right),
\end{equation}

where $M_p$ maps axes and $s_p$ sets workspace scale. The first valid device
pose becomes a reference and produces no robot movement. This prevents the
initial jump that would occur if the device pose were interpreted as an
absolute robot target.

Test one translation and one rotation axis at a time. If an axis is wrong,
correct the mapping rather than teaching the operator to compensate mentally.

\section{From target error to robot motion}

For Cartesian velocity control, a common structure is:

\begin{equation}
  v_{des} =
  \begin{bmatrix}K_p e_p & K_r e_r\end{bmatrix}^{T},
  \qquad
  \dot q = J^{\#}(q)v_{des},
\end{equation}

where $e_p$ and $e_r$ are pose errors and $J^{\#}$ is a damped pseudoinverse.
Joint velocity and acceleration limits must be applied after this conversion.

Target smoothing can reduce abrupt hand motion:

\begin{equation}
  x_f(t)=(1-\alpha)x_f(t-1)+\alpha x_{raw}(t).
\end{equation}

A large $\alpha$ is responsive but transmits more noise; a small $\alpha$ is
smooth but introduces lag. Tune scale, target-speed limit, filter alpha, servo
gain, and joint limit one at a time. Several different parameter combinations
can produce similar slow motion while having very different behavior near
contact.

\section{Understand the timing chain}

The device, controller, cameras, and recorder do not need the same frequency.
For example, the internship setup approximately used:

\begin{table}[H]
  \centering
  \caption{Typical rates in the real teleoperation pipeline.}
  \begin{tabular}{@{}ll@{}}
    \toprule
    \textbf{Process} & \textbf{Reference rate} \\
    \midrule
    Touch ROS~2 publication & approximately 100 Hz \\
    Robot control loop & 50 Hz \\
    RealSense capture & 30 Hz \\
    Demonstration recording & 10 Hz \\
    \bottomrule
  \end{tabular}
\end{table}

The controller should consume the newest available target, while the recorder
periodically samples a coherent observation-action pair. Queueing every old
device message increases lag without adding useful information.

\begin{methodbox}[Measure latency instead of guessing]
Timestamp device input, camera capture, target generation, command transmission,
and recorder write. Log dropped or stale frames. A slow response may come from
the ROS path, camera blocking, controller filtering, disk compression, or the
robot safety slider; increasing $K_p$ does not fix all of these causes.
\end{methodbox}

Useful improvements for a future implementation are:

\begin{itemize}
  \item use depth-one/latest-sample ROS QoS for interactive targets;
  \item capture cameras continuously rather than starting a blocking read in
  the servo loop;
  \item move image compression and Zarr writes to a bounded background queue;
  \item record timestamps and reject a sample when one camera is stale;
  \item display control error, frame age, and saved-step count to the operator;
  \item stop or invalidate the episode after a frozen camera rather than saving
  apparently normal repeated images.
\end{itemize}

\section{Design the gripper command as part of the task}

Continuous gripper measurements are not always good action labels. A physical
gripper moves slowly and may report different intermediate widths for the same
operator intention. If the task naturally uses semantic states, record the
commanded state as the action and the measured width as an observation.

In the internship dataset, the gripper action used three semantic states
(open, narrow, and grasp). The narrowed state allowed entry into the cavity.
The commanded state was stored in \filepath{action}, while measured opening was
stored in \filepath{robot0_gripper_qpos}; this preserved operator intention and
still exposed actuator delay.

Use continuous commands when the task genuinely requires variable grasp width.
Do not quantize only to hide inconsistent data; choose an action representation
that matches the future task.

\section{Collect demonstrations that teach observable corrections}

A good demonstration is not simply a successful one. It should make the reason
for each correction visible to the policy. Practical recommendations are:

\begin{itemize}
  \item practice the controller before recording;
  \item begin and end episodes close to the useful task interval;
  \item use a consistent phase order and start configuration;
  \item avoid moving from information that is visible only to the human;
  \item delay fine alignment until the relevant camera sees the object;
  \item avoid passing the gripper in front of the object unnecessarily;
  \item cancel failed, interrupted, or camera-corrupted demonstrations;
  \item include deliberate coverage of the expected pose range.
\end{itemize}

In the constrained-extraction experiments, early demonstrations sometimes
began orientation corrections before the wrist camera clearly saw the component. Later
demonstrations approached first, waited for a useful wrist view, and then
aligned. This made the observation-action relationship easier to learn.

\section{Launch through maintained wrappers}

The repository uses \filepath{launch_all.sh} for simulation and
\filepath{ur10_real_robot/run_teleop.sh} for physical collection. The complete
commands and controls are maintained in \filepath{README_CODE.md}. On another
platform, replace the robot adapter and limits while preserving the staged
pre-flight, recording, and dataset checks.

\begin{safetybox}[Physical teleoperation]
An operator must be trained on the platform, the speed slider must begin low,
the workspace must be clear, and an emergency stop must remain reachable. Test
each mapped axis with millimetric motions before combining directions or
recording data.
\end{safetybox}

\begin{checkpointbox}[Before a long collection]
Record three episodes, replay every frame, plot the trajectory, and verify the
gripper labels. Continue only if the motion is responsive enough for small
corrections and the saved data represents what the operator intended.
\end{checkpointbox}
